\documentclass[letterpaper,11pt]{article}
\usepackage[empty]{fullpage}
\usepackage{titlesec}
\usepackage{enumitem}
\usepackage[hidelinks]{hyperref}
\usepackage{fancyhdr}
\usepackage{fontawesome5}

\pagestyle{fancy}
\fancyhf{}
\renewcommand{\headrulewidth}{0pt}
\renewcommand{\footrulewidth}{0pt}

\addtolength{\oddsidemargin}{-0.5in}
\addtolength{\evensidemargin}{-0.5in}
\addtolength{\textwidth}{1in}
\addtolength{\topmargin}{-0.7in}
\addtolength{\textheight}{1.4in}

\titleformat{\section}
  {\scshape\raggedright\large\bfseries}
  {}{0em}{}
  [\titlerule]
\titlespacing*{\section}{0pt}{8pt}{4pt}
\setlength{\parskip}{2pt}
\setlength{\parindent}{0pt}

\begin{document}

% Choose one version to compile by uncommenting exactly one line below.
% \begin{center}
{\Huge \scshape Qiong Liu, Ph.D.} \\ \vspace{4pt}
\small
\faPhone\ 475-280-1364 \quad
\href{mailto:qiongliu.work@gmail.com}{qiongliu.work@gmail.com} \quad
\href{https://www.linkedin.com/in/qiong-l-529b23196}{LinkedIn} \quad
\href{https://scholar.google.com/citations?user=SvnpHCkAAAAJ&hl=en}{Google Scholar}

Authorized to work in the U.S. without sponsorship
\end{center}

\section{Professional Summary}
\textbf{Computer Vision / AI Researcher} with experience developing deep learning methods for image restoration, motion correction, denoising, and spatiotemporal modeling in real-world imaging systems. Strong background in \textbf{representation learning, self-supervised learning, deformable modeling, and robust evaluation} with first-author publications and patent-pending work. Experienced in translating research ideas into deployable algorithms for clinically meaningful image quality improvement.

\section{Core Expertise}
\textbf{Computer Vision \& ML:} Image restoration, denoising, deformable registration, representation learning, self-supervised learning, diffusion models, transformers \\
\textbf{Modeling:} Spatiotemporal modeling, motion modeling, signal decomposition, optimization, uncertainty-aware evaluation \\
\textbf{Programming:} Python, PyTorch, NumPy, SciPy, MATLAB \\
\textbf{Research Strengths:} Problem formulation, algorithm design, evaluation, cross-functional collaboration

\section{Experience}
\textbf{Canon Medical Research USA} \hfill Apr 2025 -- Present \\
\textit{AI Scientist (Reconstruction), Vernon Hills, IL}
\begin{itemize}[leftmargin=*]
\item Developed deep learning algorithms for \textbf{motion correction and image quality improvement} in dynamic imaging under noisy and incomplete data conditions
\item Designed \textbf{spatiotemporal modeling} methods to improve robustness and consistency across sequential image frames
\item Built training and evaluation pipelines for real-world imaging data, integrating preprocessing, model training, inference, and quantitative assessment
\item Collaborated with research and product teams to translate algorithm prototypes into deployable vision solutions
\end{itemize}

\textbf{Canon Medical Research USA} \hfill Jun 2023 -- May 2024 \\
\textit{Research Scientist Intern, Vernon Hills, IL}
\begin{itemize}[leftmargin=*]
\item Developed deep learning methods for \textbf{image denoising and low-signal restoration}, improving image fidelity while preserving critical structures
\item Designed evaluation frameworks for robustness, quantitative consistency, and generalization across heterogeneous datasets
\item Presented technical findings to guide model design and downstream product decisions
\end{itemize}

\textbf{United Imaging Healthcare} \hfill Jul 2019 -- Aug 2019 \\
\textit{Student Intern, Wuhan, China}
\begin{itemize}[leftmargin=*]
\item Supported development and validation of engineering systems for clinical imaging applications
\end{itemize}

\section{Patents}
\textbf{Edge-Guided Deformation Field Fusion for Motion Correction} \hfill 2026 \\
Patent Pending, \textit{Lead Inventor}
\begin{itemize}[leftmargin=*]
\item Proposed a spatially adaptive deformable modeling framework to improve alignment accuracy and robustness in challenging image regions
\end{itemize}

\textbf{Automatic Data-Driven Multi-Signal Decomposition} \hfill 2025 \\
Patent Pending, \textit{Lead Inventor}
\begin{itemize}[leftmargin=*]
\item Developed an AI-based framework to separate overlapping temporal signals from dynamic image sequences
\end{itemize}

\section{Selected Research Projects}

\textbf{Physiology-Guided Motion Correction for Dynamic PET}
\begin{itemize}[leftmargin=*]
\item Developed learning-based motion correction methods for dynamic PET, addressing failure modes caused by \textbf{noise disparity, mixed Gaussian/Poisson noise, and tracer-dependent uptake heterogeneity} in real patient images
\item Identified limitations of conventional network constraints based on global smoothness priors, which were often insufficient for anatomically ambiguous and low-count regions
\item Proposed a \textbf{physiology-guided mask} to provide spatial guidance during network training, improving robustness in regions with challenging tracer distribution and low image contrast
\item Introduced \textbf{spatially varying regularization} in place of a global smoothness constraint, enabling better balance between local flexibility and stable deformation estimation
\end{itemize}

\textbf{Kinetics-Guided Self-Supervised Denoising for Dynamic PET}
\begin{itemize}[leftmargin=*]
\item Developed \textbf{Patlak-Guided Self-Supervised Denoising Network (PG-SSDNet)} for dynamic PET denoising to improve the quality and quantitative reliability of parametric Ki imaging
\item Incorporated a \textbf{Patlak consistency loss} into a multi-frame self-supervised framework to preserve tracer kinetic behavior across dynamic frames
\item Demonstrated that \textbf{integrating physics/kinetic modeling into AI training} improves denoising performance while maintaining physiologically meaningful quantitative endpoints
\item Evaluated the method on whole-body \textsuperscript{18}F-FDG dynamic PET datasets and showed improved Patlak fitting consistency compared with supervised and unguided baselines
\end{itemize}

\textbf{Deep Image Prior for Dynamic PET Denoising and Quantification}
\begin{itemize}[leftmargin=*]
\item Developed population-based deep image prior methods for dynamic PET denoising, with emphasis on preserving downstream \textbf{parametric quantification accuracy}
\item Investigated how \textbf{training data composition and population priors} affect denoising behavior and quantitative outcome, showing that dataset design materially influences model performance
\item Analyzed CNN training behavior and evolving noise patterns to better understand how deep models learn image structure in low-count dynamic PET
\item Demonstrated that, although deep learning models are often treated as black boxes, their behavior can still be \textbf{systematically interpreted and leveraged} to improve both image quality and quantitative reliability
\end{itemize}

\section{Education}
\textbf{Yale University} \hfill Aug 2020 -- May 2025 \\
Ph.D. in Biomedical Engineering \\
Thesis: Improving Image Quality and Quantification Accuracy for Static and Dynamic PET

\textbf{Huazhong University of Science and Technology} \hfill Sep 2016 -- Jun 2020 \\
B.S. in Biomedical Engineering, Minor in Computer Science \\
GPA: 3.94/4.00

\section{Selected Publications}
\textbf{First Author}
\begin{itemize}[leftmargin=*]
\item Liu Q., et al., \textit{Parametric cardiac imaging with $^{18}$F-flutemetamol PET to evaluate the impact of tafamidis in patients with transthyretin cardiac amyloidosis}, \textbf{Journal of Nuclear Medicine}, 2026
\item Liu Q., et al., \textit{Patlak-guided self-supervised learning for dynamic PET denoising}, \textbf{IEEE Transactions on Radiation and Plasma Medical Sciences}, 2026
\item Liu Q., et al., \textit{Population-based deep image prior for dynamic PET denoising: a data-driven approach to improve parametric quantification}, \textbf{Medical Image Analysis}, 2024
\item Liu Q., et al., \textit{A personalized deep learning denoising strategy for low-count PET images}, \textbf{Physics in Medicine \& Biology}, 2022
\end{itemize}

\textbf{Co-Author}
\begin{itemize}[leftmargin=*]
\item Xia M., Xie H., Liu Q., et al., \textit{LeqMod: adaptable lesion-quantification-consistent modulation for deep learning low-count PET image denoising}, IEEE Transactions on Medical Imaging, 2025
\item Zhou B., Xie H., Liu Q., et al., \textit{FedFTN: personalized federated learning with deep feature transformation network for multi-institutional low-count PET denoising}, Medical Image Analysis, 2023
\end{itemize}

\section{Awards}
\textbf{Society of Nuclear Medicine and Molecular Imaging (SNMMI)}
\begin{itemize}[leftmargin=*]
\item Young Investigator Award (1st Place), 2025
\item Poster Award (2nd Place), 2024
\item Young Investigator Award (2nd Place), 2023
\end{itemize}

\textbf{Mitacs}
\begin{itemize}[leftmargin=*]
\item Globalink Research Award, 2019
\end{itemize}
% \begin{center}
{\Huge \scshape Qiong Liu, Ph.D.} \\ \vspace{4pt}
\small
\faPhone\ 475-280-1364 \quad
\href{mailto:qiongliu.work@gmail.com}{qiongliu.work@gmail.com} \quad
\href{https://www.linkedin.com/in/qiong-l-529b23196}{LinkedIn} \quad
\href{https://scholar.google.com/citations?user=SvnpHCkAAAAJ&hl=en}{Google Scholar}

Authorized to work in the U.S. without sponsorship
\end{center}

\section{Professional Summary}
\textbf{Machine Learning Engineer} with experience building end-to-end ML pipelines for image quality improvement, signal restoration, and robust modeling under noisy and incomplete data conditions. Strong background in \textbf{deep learning, self-supervised learning, evaluation design, and real-world data analysis}. Proven track record of translating technically complex problems into deployable algorithmic solutions across research and product-oriented settings.

\section{Core Expertise}
\textbf{Applied ML:} Deep learning, self-supervised learning, multimodal learning, diffusion models, transformers \\
\textbf{Image / Signal Quality:} Denoising, restoration, motion artifact reduction, quality evaluation, robustness analysis \\
\textbf{Modeling \& Evaluation:} Spatiotemporal modeling, representation learning, uncertainty modeling, quantitative validation \\
\textbf{Programming:} Python, PyTorch, NumPy, SciPy, MATLAB

\section{Experience}
\textbf{Canon Medical Research USA} \hfill Apr 2025 -- Present \\
\textit{AI Scientist (Reconstruction), Vernon Hills, IL}
\begin{itemize}[leftmargin=*]
\item Built machine learning pipelines for \textbf{image quality improvement} and motion-related degradation correction in real-world imaging data
\item Developed robust algorithms for noisy, low-signal, and incomplete data settings, improving downstream consistency and usability
\item Designed quantitative evaluation workflows to assess model performance, generalization, and reliability across data conditions
\item Worked cross-functionally to move algorithm development toward production-oriented deployment
\end{itemize}

\textbf{Canon Medical Research USA} \hfill Jun 2023 -- May 2024 \\
\textit{Research Scientist Intern, Vernon Hills, IL}
\begin{itemize}[leftmargin=*]
\item Developed deep learning models for \textbf{signal restoration and denoising}, preserving critical structures under challenging acquisition conditions
\item Designed validation frameworks for robustness, consistency, and quality assessment across heterogeneous datasets
\item Communicated technical results to support algorithm design and product decisions
\end{itemize}

\textbf{United Imaging Healthcare} \hfill Jul 2019 -- Aug 2019 \\
\textit{Student Intern, Wuhan, China}
\begin{itemize}[leftmargin=*]
\item Supported development and validation of engineering systems for clinical applications
\end{itemize}

\section{Patents}
\textbf{Edge-Guided Deformation Field Fusion for Motion Correction} \hfill 2026 \\
Patent Pending, \textit{Lead Inventor}
\begin{itemize}[leftmargin=*]
\item Developed a spatially adaptive algorithm to improve robustness and alignment quality in complex image regions
\end{itemize}

\textbf{Automatic Data-Driven Multi-Signal Decomposition} \hfill 2025 \\
Patent Pending, \textit{Lead Inventor}
\begin{itemize}[leftmargin=*]
\item Developed an ML-based method to separate overlapping temporal signals from unlabeled dynamic data
\end{itemize}

\section{Selected Machine Learning Projects}

\textbf{Robust Motion-Aware Image Modeling}
\begin{itemize}[leftmargin=*]
\item Developed ML methods for motion-related image degradation correction under noisy, low-signal, and heterogeneous real-world conditions
\item Proposed spatially adaptive regularization and structured spatial guidance to improve model robustness beyond conventional globally constrained approaches
\item Improved consistency and reliability of downstream image analysis in challenging temporal data
\end{itemize}

\textbf{Self-Supervised Temporal Denoising}
\begin{itemize}[leftmargin=*]
\item Built self-supervised denoising models for temporal imaging data with explicit consistency constraints across frames
\item Incorporated domain-informed priors into training to improve robustness, generalization, and preservation of meaningful temporal structure
\item Demonstrated that structured constraints can improve model performance in noisy and limited-data regimes
\end{itemize}

\textbf{Interpretable Deep Learning for Restoration}
\begin{itemize}[leftmargin=*]
\item Developed CNN-based and deep image prior models for denoising and reconstruction across heterogeneous data conditions
\item Studied the effects of training data composition and optimization behavior on restoration quality and model reliability
\item Analyzed model behavior during training to better understand how deep networks learn signal and noise structure
\end{itemize}

\section{Education}
\textbf{Yale University} \hfill Aug 2020 -- May 2025 \\
Ph.D. in Biomedical Engineering \\
Thesis: Improving Image Quality and Quantification Accuracy for Static and Dynamic PET

\textbf{Huazhong University of Science and Technology} \hfill Sep 2016 -- Jun 2020 \\
B.S. in Biomedical Engineering, Minor in Computer Science \\
GPA: 3.94/4.00

\section{Selected Publications}
\begin{itemize}[leftmargin=*]
\item Liu Q., et al., \textit{Patlak-guided self-supervised learning for dynamic PET denoising}, \textbf{IEEE Transactions on Radiation and Plasma Medical Sciences}, 2026
\item Liu Q., et al., \textit{Population-based deep image prior for dynamic PET denoising: a data-driven approach to improve parametric quantification}, \textbf{Medical Image Analysis}, 2024
\item Liu Q., et al., \textit{A personalized deep learning denoising strategy for low-count PET images}, \textbf{Physics in Medicine \& Biology}, 2022
\end{itemize}

\section{Awards}
\textbf{Society of Nuclear Medicine and Molecular Imaging (SNMMI)}
\begin{itemize}[leftmargin=*]
\item Young Investigator Award (1st Place), 2025
\item Poster Award (2nd Place), 2024
\item Young Investigator Award (2nd Place), 2023
\end{itemize}

\textbf{Mitacs}
\begin{itemize}[leftmargin=*]
\item Globalink Research Award, 2019
\end{itemize}
% \begin{center}
{\Huge \scshape Qiong Liu, Ph.D.} \\ \vspace{4pt}
\small
\faPhone\ 475-280-1364 \quad
\href{mailto:qiongliu.work@gmail.com}{qiongliu.work@gmail.com} \quad
\href{https://www.linkedin.com/in/qiong-l-529b23196}{LinkedIn} \quad
\href{https://scholar.google.com/citations?user=SvnpHCkAAAAJ&hl=en}{Google Scholar}

Authorized to work in the U.S. without sponsorship
\end{center}

\section{Professional Summary}
\textbf{Machine Learning / Algorithm Engineer} with experience developing data-driven models for \textbf{noisy, high-dimensional, and temporally structured data}. Strong background in \textbf{signal processing, self-supervised learning, anomaly detection, and spatiotemporal modeling} for real-world systems. Experienced in building robust pipelines for \textbf{feature extraction, model training, and quantitative evaluation} under limited and imperfect data conditions. Skilled at translating complex data patterns into actionable insights for system optimization.

\section{Core Expertise}
\textbf{Machine Learning:} Representation learning, self-supervised learning, anomaly detection, contrastive learning \\
\textbf{Signal \& Data Modeling:} Time-series analysis, spatiotemporal modeling, signal decomposition, noise modeling \\
\textbf{Computer Vision:} Image restoration, defect detection, pattern recognition, feature extraction \\
\textbf{Programming:} Python, PyTorch, NumPy, SciPy, MATLAB \\
\textbf{Engineering Strengths:} Model robustness, generalization, data-driven system optimization

\section{Experience}

\textbf{Canon Medical Research USA} \hfill Apr 2025 -- Present \\
\textit{Reconstruction Scientist, Vernon Hills, IL}
\begin{itemize}[leftmargin=*]
\item Developed methods to \textbf{extract dual motion signals} from noisy dynamic data with overlapping temporal patterns and distribution shift
\item Built \textbf{signal-aware motion correction models} to address noise--feature trade-offs using domain-informed constraints
\item Developed \textbf{diffusion-based denoising models} and integrated them into reconstruction pipelines to improve overall data quality
\end{itemize}

\textbf{Canon Medical Research USA} \hfill Jun 2023 -- May 2024 \\
\textit{Research Scientist Intern, Vernon Hills, IL}
\begin{itemize}[leftmargin=*]
\item Developed \textbf{efficient attention-based denoising models} with modified transformer architectures for faster inference
\item Designed \textbf{physics-informed noise modeling methods} for signal-dependent noise in low-SNR data
\item Built deep learning models for \textbf{scatter estimation}, reducing a major runtime bottleneck in the reconstruction pipeline
\end{itemize}

\textbf{Yale University, Yale PET Center} \hfill Aug 2020 -- May 2025 \\
\textit{Ph.D. Researcher (Machine Learning / Imaging Systems), New Haven, CT}
\begin{itemize}[leftmargin=*]
\item Developed learning-based methods for \textbf{low-SNR signal reconstruction and denoising}, using self-supervised learning and domain-informed constraints to preserve meaningful structure
\item Built robust learning pipelines for \textbf{data-scarce and heterogeneous settings}, including federated learning, fine-tuning, and adaptation across varying data distributions
\item Developed \textbf{dynamic system modeling} methods for temporal signal decomposition, quantitative analysis, and data-driven modeling of novel tracer behavior
\end{itemize}

\textbf{United Imaging Healthcare} \hfill May 2019 -- Aug 2019 \\
\textit{Student Intern, Wuhan, China}
\begin{itemize}[leftmargin=*]
\item Supported development and testing of a \textbf{surgical robotics system}, including image registration, system validation, and performance evaluation
\end{itemize}

\section{Patents}
\textbf{Edge-Guided Deformation Field Fusion for Motion Correction} \hfill 2026 \\
Patent Pending, \textit{Lead Inventor}
\begin{itemize}[leftmargin=*]
\item Proposed a spatially adaptive deformable modeling framework to improve alignment accuracy and robustness in challenging image regions
\end{itemize}

\textbf{Automatic Data-Driven Multi-Signal Decomposition} \hfill 2025 \\
Patent Pending, \textit{Lead Inventor}
\begin{itemize}[leftmargin=*]
\item Developed an AI-based framework to separate overlapping temporal signals from dynamic image sequences
\end{itemize}

\section{Selected Projects}

\textbf{Self-Supervised Denoising with Physics-Guided Constraints}
\begin{itemize}[leftmargin=*]
\item Developed a self-supervised framework for \textbf{low-SNR signal denoising} using multi-frame consistency and reconstruction objectives
\item Introduced \textbf{physics-guided constraints} (e.g., kinetic consistency) to preserve meaningful signal behavior during denoising
\item Demonstrated improved robustness and quantitative consistency compared to supervised and unguided methods
\end{itemize}

\textbf{Data-Driven Prior Modeling for Low-Signal Reconstruction}
\begin{itemize}[leftmargin=*]
\item Developed population-based and data-driven priors for \textbf{signal reconstruction under extreme noise and limited data}
\item Analyzed how \textbf{training data distribution and prior design} impact model behavior and reconstruction quality
\item Improved reconstruction fidelity while maintaining downstream quantitative accuracy
\end{itemize}

\textbf{Dynamic System Modeling for Quantitative Signal Analysis}
\begin{itemize}[leftmargin=*]
\item Developed \textbf{parametric modeling and system identification methods} for dynamic processes in sequential data
\item Applied modeling techniques to extract \textbf{quantitative features} for analysis and evaluation of complex systems
\item Enabled improved interpretability and decision-making through data-driven modeling of dynamic behavior
\end{itemize}

\section{Education}
\textbf{Yale University} \hfill Aug 2020 -- May 2025 \\
Ph.D. in Biomedical Engineering \\
Thesis: Improving Image Quality and Quantification Accuracy for Static and Dynamic PET

\textbf{Huazhong University of Science and Technology} \hfill Sep 2016 -- Jun 2020 \\
B.S. in Biomedical Engineering, Minor in Computer Science \\
GPA: 3.94/4.00

\section{Selected Publications}

\textbf{First Author}
\begin{itemize}[leftmargin=*]
\item Liu Q., et al., \textit{Parametric cardiac imaging with $^{18}$F-flutemetamol PET to evaluate the impact of tafamidis in patients with transthyretin cardiac amyloidosis}, \textbf{Journal of Nuclear Medicine}, 2026
\item Liu Q., et al., \textit{Patlak-guided self-supervised learning for dynamic PET denoising}, \textbf{IEEE Transactions on Radiation and Plasma Medical Sciences}, 2026
\item Liu Q., et al., \textit{Population-based deep image prior for dynamic PET denoising: a data-driven approach to improve parametric quantification}, \textbf{Medical Image Analysis}, 2024
\item Liu Q., et al., \textit{A personalized deep learning denoising strategy for low-count PET images}, \textbf{Physics in Medicine \& Biology}, 2022
\end{itemize}

\textbf{Co-Author}
\begin{itemize}[leftmargin=*]
\item Xia M., Xie H., Liu Q., et al., \textit{LeqMod: adaptable lesion-quantification-consistent modulation for deep learning low-count PET image denoising}, IEEE Transactions on Medical Imaging, 2025
\item Zhou B., Xie H., Liu Q., et al., \textit{FedFTN: personalized federated learning with deep feature transformation network for multi-institutional low-count PET denoising}, Medical Image Analysis, 2023
\end{itemize}

\section{Awards}
\textbf{Society of Nuclear Medicine and Molecular Imaging (SNMMI)}
\begin{itemize}[leftmargin=*]
\item Young Investigator Award (1st Place), 2025
\item Poster Award (2nd Place), 2024
\item Young Investigator Award (2nd Place), 2023
\end{itemize}

\textbf{Mitacs}
\begin{itemize}[leftmargin=*]
\item Globalink Research Award, 2019
\end{itemize}
% \begin{center}
{\Huge \scshape Qiong Liu, Ph.D.} \\ \vspace{4pt}
\small
\faPhone\ 475-280-1364 \quad
\href{mailto:qiongliu.work@gmail.com}{qiongliu.work@gmail.com} \quad
\href{https://www.linkedin.com/in/qiong-l-529b23196}{LinkedIn} \quad
\href{https://scholar.google.com/citations?user=SvnpHCkAAAAJ&hl=en}{Google Scholar}

Authorized to work in the U.S. without sponsorship
\end{center}

\section{Professional Summary}
\textbf{Healthcare AI / Medical Imaging Scientist} with expertise in machine learning for image reconstruction, denoising, motion correction, and quantitative analysis in PET imaging. Strong track record of developing clinically meaningful AI methods for low-count, dynamic, and motion-corrupted imaging data. Experienced across research and product-oriented environments, with first-author publications, patent-pending work, and end-to-end algorithm development experience.

\section{Core Expertise}
\textbf{Medical Imaging AI:} PET reconstruction, denoising, motion correction, quantitative imaging, parametric imaging, dynamic imaging \\
\textbf{Machine Learning:} Deep learning, self-supervised learning, diffusion models, transformers, multimodal learning, physics-informed AI \\
\textbf{Programming:} Python, PyTorch, NumPy, SciPy, MATLAB \\
\textbf{Strengths:} Algorithm development, quantitative evaluation, clinically informed modeling, cross-functional collaboration

\section{Experience}
\textbf{Canon Medical Research USA} \hfill Apr 2025 -- Present \\
\textit{AI Scientist (Reconstruction), Vernon Hills, IL}
\begin{itemize}[leftmargin=*]
\item Developing AI-driven motion correction methods for PET reconstruction to improve image quality and quantitative accuracy in dynamic and low-count imaging
\item Designing next-generation gating algorithms for respiratory and cardiac motion compensation in dynamic PET
\item Building reconstruction pipelines that integrate deep learning and physics-based modeling for clinically robust imaging
\item Collaborating across research and product teams to evaluate algorithm performance and translational potential
\end{itemize}

\textbf{Canon Medical Research USA} \hfill Jun 2023 -- May 2024 \\
\textit{Research Scientist Intern, Vernon Hills, IL}
\begin{itemize}[leftmargin=*]
\item Developed AI-based PET reconstruction and denoising methods for low-count imaging scenarios
\item Designed evaluation frameworks to assess image quality and quantitative consistency across datasets
\item Presented technical findings to guide product and algorithm development
\end{itemize}

\textbf{United Imaging Healthcare} \hfill Jul 2019 -- Aug 2019 \\
\textit{Student Intern, Wuhan, China}
\begin{itemize}[leftmargin=*]
\item Supported development and validation of engineering systems for clinical applications
\end{itemize}

\section{Patents}
\textbf{Edge-Guided Deformation Field Fusion for PET Motion Correction} \hfill 2026 \\
Patent Pending, \textit{Lead Inventor}
\begin{itemize}[leftmargin=*]
\item Proposed a spatially adaptive AI framework to improve lesion localization and quantitative accuracy in PET motion correction
\end{itemize}

\textbf{Automatic Data-Driven Cardiac-Respiratory Gating} \hfill 2025 \\
Patent Pending, \textit{Lead Inventor}
\begin{itemize}[leftmargin=*]
\item Developed an AI-based method for separating respiratory and cardiac motion signals from dynamic PET data
\end{itemize}

\section{Selected Projects}

\textbf{Physiology-Guided Motion Correction for Dynamic PET}
\begin{itemize}[leftmargin=*]
\item Developed AI-based motion correction methods for dynamic PET, addressing real-world challenges from heterogeneous tracer uptake, low-count noise, and anatomically ambiguous regions
\item Proposed a \textbf{physiology-guided spatial mask} to provide clinically meaningful guidance during network training, improving robustness in regions where conventional learning-based constraints were insufficient
\item Introduced \textbf{spatially adaptive regularization} in place of global smoothness constraints, enabling more reliable deformation estimation across heterogeneous image regions
\item Improved motion correction robustness and temporal consistency for downstream quantitative analysis in dynamic PET imaging
\end{itemize}

\textbf{Kinetics-Guided Dynamic PET Denoising}
\begin{itemize}[leftmargin=*]
\item Developed \textbf{Patlak-guided self-supervised learning} methods for dynamic PET denoising to improve the quality and quantitative reliability of parametric imaging
\item Incorporated tracer kinetic constraints directly into AI model training, preserving physiologically meaningful temporal behavior across dynamic frames
\item Demonstrated that integrating \textbf{physics- and kinetics-informed priors} into deep learning improves performance in noisy and limited-data clinical imaging settings
\item Improved downstream Patlak-based quantification and parametric consistency in whole-body dynamic \textsuperscript{18}F-FDG PET studies
\end{itemize}

\textbf{Interpretable Deep Learning for Low-Count and Dynamic PET}
\begin{itemize}[leftmargin=*]
\item Developed personalized, population-based, and deep image prior methods for PET denoising with emphasis on preserving clinically relevant quantitative information
\item Studied how training data composition and population priors affect denoising outcome and downstream quantification performance across heterogeneous imaging datasets
\item Analyzed CNN training behavior and evolving noise patterns to better understand model behavior in low-count and dynamic PET reconstruction tasks
\item Showed that deep learning for healthcare imaging, while often treated as a black box, can be systematically interpreted and optimized for more reliable clinical use
\end{itemize}

\section{Education}
\textbf{Yale University} \hfill Aug 2020 -- May 2025 \\
Ph.D. in Biomedical Engineering \\
Thesis: Improving Image Quality and Quantification Accuracy for Static and Dynamic PET

\textbf{Huazhong University of Science and Technology} \hfill Sep 2016 -- Jun 2020 \\
B.S. in Biomedical Engineering, Minor in Computer Science \\
GPA: 3.94/4.00

\section{Selected Publications}
\textbf{First Author}
\begin{itemize}[leftmargin=*]
\item Liu Q., et al., \textit{Parametric cardiac imaging with $^{18}$F-flutemetamol PET to evaluate the impact of tafamidis in patients with transthyretin cardiac amyloidosis}, \textbf{Journal of Nuclear Medicine}, 2026
\item Liu Q., et al., \textit{Patlak-guided self-supervised learning for dynamic PET denoising}, \textbf{IEEE Transactions on Radiation and Plasma Medical Sciences}, 2026
\item Liu Q., et al., \textit{Population-based deep image prior for dynamic PET denoising: a data-driven approach to improve parametric quantification}, \textbf{Medical Image Analysis}, 2024
\item Liu Q., et al., \textit{A personalized deep learning denoising strategy for low-count PET images}, \textbf{Physics in Medicine \& Biology}, 2022
\end{itemize}

\section{Awards}
\textbf{Society of Nuclear Medicine and Molecular Imaging (SNMMI)}
\begin{itemize}[leftmargin=*]
\item Young Investigator Award (1st Place), 2025
\item Poster Award (2nd Place), 2024
\item Young Investigator Award (2nd Place), 2023
\end{itemize}

\textbf{Mitacs}
\begin{itemize}[leftmargin=*]
\item Globalink Research Award, 2019
\end{itemize}
% \begin{center}
{\Huge \scshape Qiong Liu, Ph.D.} \\ \vspace{4pt}
\small
\faPhone\ 475-280-1364 \quad
\href{mailto:qiongliu.work@gmail.com}{qiongliu.work@gmail.com} \quad
\href{https://www.linkedin.com/in/qiong-l-529b23196}{LinkedIn} \quad
\href{https://scholar.google.com/citations?user=SvnpHCkAAAAJ&hl=en}{Google Scholar}

Authorized to work in the U.S. without sponsorship
\end{center}

\section{Professional Summary}
\textbf{Machine Learning Researcher} with experience developing learning-based methods for spatiotemporal modeling, representation learning, signal decomposition, and robust inference from noisy real-world data. Strong background in deep learning, self-supervised learning, and data-efficient modeling, with first-author publications and patent-pending work. Experienced in taking research ideas from problem formulation through empirical validation and product-oriented collaboration.

\section{Core Expertise}
\textbf{Machine Learning:} Deep learning, self-supervised learning, multimodal learning, transformers, diffusion models, representation learning \\
\textbf{Research Areas:} Spatiotemporal modeling, temporal representation learning, signal decomposition, optimization, uncertainty modeling \\
\textbf{Programming:} Python, PyTorch, NumPy, SciPy, MATLAB \\
\textbf{Strengths:} Algorithm design, research experimentation, evaluation, cross-functional collaboration

\section{Experience}
\textbf{Canon Medical Research USA} \hfill Apr 2025 -- Present \\
\textit{AI Scientist (Reconstruction), Vernon Hills, IL}
\begin{itemize}[leftmargin=*]
\item Led development of learning-based methods for \textbf{spatiotemporal modeling} under noisy and incomplete data conditions
\item Designed \textbf{multi-signal decomposition} frameworks to separate overlapping temporal patterns for downstream modeling
\item Built end-to-end ML pipelines integrating data processing, model training, inference, and evaluation
\item Collaborated across research and product teams to translate algorithmic ideas into deployable systems
\end{itemize}

\textbf{Canon Medical Research USA} \hfill Jun 2023 -- May 2024 \\
\textit{Research Scientist Intern, Vernon Hills, IL}
\begin{itemize}[leftmargin=*]
\item Developed data-efficient deep learning models for low-signal regimes, improving robustness while preserving critical information
\item Designed evaluation frameworks to assess generalization, consistency, and quantitative reliability across datasets
\item Delivered technical findings to guide algorithm design and applied research directions
\end{itemize}

\textbf{United Imaging Healthcare} \hfill Jul 2019 -- Aug 2019 \\
\textit{Student Intern, Wuhan, China}
\begin{itemize}[leftmargin=*]
\item Supported development and validation of engineering systems for clinical applications
\end{itemize}

\section{Patents}
\textbf{Edge-Guided Deformation Field Fusion for Motion Correction} \hfill 2026 \\
Patent Pending, \textit{Lead Inventor}
\begin{itemize}[leftmargin=*]
\item Proposed a spatially adaptive learning framework to improve alignment robustness in complex dynamic data
\end{itemize}

\textbf{Automatic Data-Driven Multi-Signal Decomposition} \hfill 2025 \\
Patent Pending, \textit{Lead Inventor}
\begin{itemize}[leftmargin=*]
\item Developed a representation learning approach for separating overlapping temporal signals from unlabeled data
\end{itemize}

\section{Selected Research Projects}

\textbf{Structured Learning for Spatiotemporal Alignment}
\begin{itemize}[leftmargin=*]
\item Developed learning-based methods for alignment across sequential high-dimensional observations with challenging noise and correspondence ambiguity
\item Proposed \textbf{spatially adaptive regularization} to replace globally uniform constraints, improving robustness in low-signal and ambiguous regions
\item Studied how \textbf{structured spatial guidance} can improve optimization and generalization when standard learning objectives fail under real-world variability
\item Designed pairwise and multi-frame learning strategies for scalable temporal correspondence modeling
\end{itemize}

\textbf{Self-Supervised Learning with Physics-Informed Constraints}
\begin{itemize}[leftmargin=*]
\item Developed self-supervised temporal denoising methods that preserve meaningful latent dynamics across sequential observations
\item Incorporated \textbf{structured consistency constraints} into learning objectives, showing that physics- and domain-informed priors can substantially improve robustness in noisy, limited-data settings
\item Demonstrated improved stability and downstream reliability compared with purely data-driven baselines
\end{itemize}

\textbf{Interpretable Data-Efficient Restoration and Representation Learning}
\begin{itemize}[leftmargin=*]
\item Developed CNN-based and deep image prior approaches for restoration under heterogeneous data distributions and low-signal regimes
\item Investigated the role of \textbf{inductive bias, training data composition, and optimization dynamics} in determining model behavior and output quality
\item Analyzed model training trajectories and evolving noise structure to better understand how deep networks learn signal representations in underconstrained problems
\end{itemize}

\section{Education}
\textbf{Yale University} \hfill Aug 2020 -- May 2025 \\
Ph.D. in Biomedical Engineering \\
Thesis: Improving Image Quality and Quantification Accuracy for Static and Dynamic PET

\textbf{Huazhong University of Science and Technology} \hfill Sep 2016 -- Jun 2020 \\
B.S. in Biomedical Engineering, Minor in Computer Science \\
GPA: 3.94/4.00

\section{Selected Publications}
\textbf{First Author}
\begin{itemize}[leftmargin=*]
\item Liu Q., et al., \textit{Patlak-guided self-supervised learning for dynamic PET denoising}, \textbf{IEEE Transactions on Radiation and Plasma Medical Sciences}, 2026
\item Liu Q., et al., \textit{Population-based deep image prior for dynamic PET denoising: a data-driven approach to improve parametric quantification}, \textbf{Medical Image Analysis}, 2024
\item Liu Q., et al., \textit{A personalized deep learning denoising strategy for low-count PET images}, \textbf{Physics in Medicine \& Biology}, 2022
\end{itemize}

\section{Awards}
\textbf{Society of Nuclear Medicine and Molecular Imaging (SNMMI)}
\begin{itemize}[leftmargin=*]
\item Young Investigator Award (1st Place), 2025
\item Poster Award (2nd Place), 2024
\item Young Investigator Award (2nd Place), 2023
\end{itemize}

\textbf{Mitacs}
\begin{itemize}[leftmargin=*]
\item Globalink Research Award, 2019
\end{itemize}
\begin{center}
{\Huge \scshape Qiong Liu, Ph.D.} \\ \vspace{4pt}
\small
\faPhone\ 475-280-1364 \quad
\href{mailto:qiongliu.work@gmail.com}{qiongliu.work@gmail.com} \quad
\href{https://www.linkedin.com/in/qiong-l-529b23196}{LinkedIn} \quad
\href{https://scholar.google.com/citations?user=SvnpHCkAAAAJ&hl=en}{Google Scholar}

Authorized to work in the U.S. without sponsorship
\end{center}

\section{Professional Summary}
\textbf{AI Scientist and Biomedical Imaging Researcher} with expertise in \textbf{cardiovascular AI, cardiac PET imaging, respiratory/cardiac motion correction, physiological signal extraction, tracer kinetic modeling, and quantitative imaging analysis}. Experienced in developing deep learning and physics-informed methods for \textbf{cardiac motion compensation, low-count denoising, parametric imaging, and clinically robust quantitative analysis} in dynamic PET. Strong background in translating cardiovascular imaging workflows into AI-driven solutions through close collaboration with cardiologists, physicians, imaging scientists, and product teams.

\section{Core Expertise}
\textbf{Cardiovascular AI \& Quantitative Imaging:} Cardiac PET imaging, tracer kinetic modeling, parametric imaging, physiological signal analysis, ECG-correlated validation, quantitative biomarker extraction \\
\textbf{Motion Correction \& Signal Processing:} Respiratory/cardiac dual gating, frequency-domain analysis, dynamic PET motion compensation, deformation field fusion \\
\textbf{Machine Learning:} Deep learning, self-supervised learning, physics-informed AI, representation learning, diffusion models, transformers \\
\textbf{Medical Imaging:} PET reconstruction, low-count denoising, dynamic PET, image quality assessment, clinical imaging validation \\
\textbf{Programming:} Python, PyTorch, NumPy, SciPy, MATLAB

\section{Experience}
\textbf{Canon Medical Research USA} \hfill Apr 2025 -- Present \\
\textit{AI Scientist / Reconstruction Scientist, Vernon Hills, IL}
\begin{itemize}[leftmargin=*]
\item Developed AI-driven motion correction and reconstruction methods for \textbf{dynamic cardiac PET imaging} under noisy and low-count acquisition conditions.
\item Designed \textbf{automatic cardiac-respiratory dual gating algorithms} to extract physiological motion signals directly from dynamic PET sequences.
\item Developed \textbf{EdgeFuse}, an edge-guided deformation field fusion framework for cardiac PET motion correction to improve myocardial boundary preservation and quantitative accuracy.
\item Built end-to-end pipelines integrating preprocessing, model training, inference, motion correction, denoising, and downstream quantitative assessment.
\item Collaborated with physicians, imaging scientists, and product teams to evaluate image quality, temporal consistency, quantitative robustness, and clinical usability.
\end{itemize}

\textbf{Canon Medical Research USA} \hfill Jun 2023 -- May 2024 \\
\textit{Research Scientist Intern, Vernon Hills, IL}
\begin{itemize}[leftmargin=*]
\item Developed deep learning methods for low-count PET denoising while preserving quantitative fidelity for downstream dynamic and parametric analysis.
\item Designed evaluation frameworks to assess robustness, generalization, image quality, and quantitative consistency across heterogeneous imaging datasets.
\item Presented technical findings to guide algorithm design and translational imaging research.
\end{itemize}

\textbf{Yale University PET Center} \hfill Aug 2020 -- May 2025 \\
\textit{Ph.D. Researcher, New Haven, CT}
\begin{itemize}[leftmargin=*]
\item Developed AI and quantitative imaging methods for \textbf{dynamic cardiac PET analysis, tracer kinetic modeling, parametric imaging, and low-count denoising}.
\item Applied \textbf{Patlak analysis} and kinetics-informed methods to derive image-based cardiovascular biomarkers from noisy dynamic PET data.
\item Performed quantitative cardiac PET analysis to evaluate treatment-related effects in transthyretin cardiac amyloidosis.
\item Built self-supervised and deep image prior-based denoising methods that preserve kinetic modeling accuracy and downstream parametric quantification.
\item Collaborated with cardiologists, nuclear medicine physicians, and clinical researchers to define clinically meaningful imaging endpoints and validate quantitative results.
\end{itemize}

\textbf{United Imaging Healthcare} \hfill Jul 2019 -- Aug 2019 \\
\textit{Student Intern, Wuhan, China}
\begin{itemize}[leftmargin=*]
\item Supported development and validation of engineering systems for clinical applications.
\end{itemize}

\section{Selected Cardiovascular AI Projects}

\textbf{Automatic Cardiac-Respiratory Dual Gating for Dynamic PET}
\begin{itemize}[leftmargin=*]
\item Developed an AI-driven dual gating framework for simultaneous extraction of respiratory and cardiac motion signals directly from dynamic PET data.
\item Combined autoencoder feature learning, PCA, and frequency-domain analysis to separate physiological motion components.
\item Validated extracted cardiac signals against ECG recordings and visually evaluated respiratory motion using high-temporal-resolution PET mini-frame movies.
\end{itemize}

\textbf{EdgeFuse: Edge-Guided Motion Correction for Cardiac PET}
\begin{itemize}[leftmargin=*]
\item Developed a spatially adaptive deep learning framework for cardiac PET motion correction using edge-guided deformation field fusion.
\item Combined deformation fields with different smoothness constraints to preserve myocardial boundaries while maintaining smooth background motion.
\item Improved motion compensation, lesion localization, and quantitative accuracy in cardiac PET imaging.
\end{itemize}

\textbf{Dynamic Cardiac PET Quantification and Kinetic Modeling}
\begin{itemize}[leftmargin=*]
\item Developed quantitative analysis pipelines for dynamic cardiac PET using Patlak-guided and kinetics-informed methods.
\item Evaluated treatment-related effects in transthyretin cardiac amyloidosis using parametric imaging biomarkers.
\item Improved robustness of quantitative cardiovascular imaging under noisy and low-count acquisition conditions.
\end{itemize}

\textbf{Low-Count and Dynamic PET Denoising}
\begin{itemize}[leftmargin=*]
\item Developed personalized, population-based, and self-supervised denoising frameworks for dynamic and low-count PET imaging.
\item Preserved quantitative fidelity and downstream kinetic modeling accuracy during denoising.
\item Designed evaluation protocols emphasizing image quality, temporal consistency, and parametric quantification.
\end{itemize}

\section{Patents}
\textbf{Edge-Guided Deformation Field Fusion for PET Motion Correction} \hfill Patent Pending \\
\textit{Lead Inventor}
\begin{itemize}[leftmargin=*]
\item Developed an AI-driven cardiac PET motion correction framework that adaptively combines deformation fields to preserve myocardial boundaries and improve quantitative accuracy.
\end{itemize}

\textbf{Automatic Data-Driven Cardiac-Respiratory Gating} \hfill Patent Pending \\
\textit{Lead Inventor}
\begin{itemize}[leftmargin=*]
\item Developed an AI-based dual gating framework for simultaneous extraction of respiratory and cardiac physiological motion signals directly from dynamic PET data.
\end{itemize}

\section{Education}
\textbf{Yale University} \hfill Aug 2020 -- May 2025 \\
Ph.D. in Biomedical Engineering \\
Thesis: Improving Image Quality and Quantification Accuracy for Static and Dynamic PET

\textbf{Huazhong University of Science and Technology} \hfill Sep 2016 -- Jun 2020 \\
B.S. in Biomedical Engineering, Minor in Computer Science \\
GPA: 3.94/4.00

\section{Selected Publications}
\begin{itemize}[leftmargin=*]
\item Liu Q., et al., \textit{Parametric cardiac imaging with $^{18}$F-flutemetamol PET to evaluate the impact of tafamidis in patients with transthyretin cardiac amyloidosis}, \textbf{Journal of Nuclear Medicine}, 2026.
\item Liu Q., et al., \textit{Patlak-guided self-supervised learning for dynamic PET denoising}, \textbf{IEEE Transactions on Radiation and Plasma Medical Sciences}, 2026.
\item Liu Q., et al., \textit{Population-based deep image prior for dynamic PET denoising: a data-driven approach to improve parametric quantification}, \textbf{Medical Image Analysis}, 2024.
\item Liu Q., et al., \textit{A personalized deep learning denoising strategy for low-count PET images}, \textbf{Physics in Medicine \& Biology}, 2022.
\item Liu Q., et al., \textit{Dynamic imaging and tracer kinetic modeling of $^{18}$F-flutemetamol PET for ATTR cardiac amyloidosis patients}, \textbf{SNMMI Oral Presentation}, 2023.
\item Liu Q., et al., \textit{Comparison of network structures for deep learning-based cardiac PET scatter correction}, \textbf{IEEE NSS MIC}, 2024.
\end{itemize}

\section{Awards}
\textbf{Society of Nuclear Medicine and Molecular Imaging (SNMMI)}
\begin{itemize}[leftmargin=*]
\item Young Investigator Award, 1st Place, Cardiovascular Council Clinical Science, 2025.
\item Poster Award, 2nd Place, Physics, Instrumentation \& Data Sciences Council, 2024.
\item Young Investigator Award, 2nd Place, Cardiovascular Council Clinical Science, 2023.
\end{itemize}

\textbf{Mitacs}
\begin{itemize}[leftmargin=*]
\item Globalink Research Award, 2019.
\end{itemize}

\end{document}
% I'm targeting a base salary in the $160K–$180K range, depending on the full compensation package.

% Canon Medical Research USA — Apr 2025 to Present
Machine Learning Scientist / Algorithm Engineer, Vernon Hills, IL

Developed methods to extract dual motion signals from noisy dynamic data with overlapping temporal patterns and distribution shift. Built signal-aware motion correction models to address noise–feature trade-offs using domain-informed constraints. Developed diffusion-based denoising models and integrated them into reconstruction pipelines to improve overall data quality

Canon Medical Research USA — Jun 2023 to May 2024
Research Scientist Intern, Vernon Hills, IL

Developed efficient attention-based denoising models with modified transformer architectures for faster inference. Designed physics-informed noise modeling methods for signal-dependent noise in low-SNR data. Built deep learning models for scatter estimation, reducing a major runtime bottleneck in the reconstruction pipeline.

Supported development and testing of a surgical robotics system, including image registration, system validation, and performance evaluation

% Hi, I’m Qiong Liu. I received my PhD from Yale University and am currently working as an AI Scientist at Canon Medical Research USA. My background is in machine learning, with experience building end-to-end pipelines for image quality improvement, signal restoration, and robust modeling under noisy and incomplete real-world data conditions. I’ve worked on deep learning, self-supervised learning, and evaluation design across both research and product-oriented settings, and I’m especially interested in developing practical ML solutions for challenging real-world problems.

%   Developed learning-based methods for low-SNR signal reconstruction and denoising, using self-supervised learning and domain-informed constraints to preserve meaningful structure. Built robust learning pipelines for data-scarce and heterogeneous settings, including federated learning, fine-tuning, and adaptation across varying data distributions. Developed dynamic system modeling methods for temporal signal decomposition, quantitative analysis, and data-driven modeling of novel tracer behavior


Supported development and testing of a surgical robotics system, including image registration, system validation, and performance evaluation.